\documentclass[10pt,a4paper]{amsart}
\usepackage{amsmath,amssymb,amsthm,graphicx}
% Retained minimal preamble for the Stacks Project chapter more-algebra.tex (More on Algebra).
% The chapter file pinned-original-source.tex carries only \input{preamble} on its first line, so no
% preamble text lives inside the chapter itself.  The Stacks Project uses one single shared preamble for
% every chapter; the snapshot stages it identically as sources/stacks-chow/preamble.tex and
% sources/stacks-etale-cohomology/preamble.tex.  The block below is copied mechanically, line for line,
% from sources/stacks-chow/preamble.tex lines 14-18 (the Xy-pic setup) and 154-217 (the native theorem
% declarations together with the mathematical macros).  Nothing was added, removed or redefined; the
% wrapper class and the amsmath/amssymb/amsthm/graphicx base packages are supplied by the assembler.
% Chapter: more-algebra.tex -- More on Algebra.
% For commutative diagrams we use Xy-pic
\usepackage[all]{xy}

% We use 2cell for 2-commutative diagrams.
\xyoption{2cell}

% Theorem environments.
%
\theoremstyle{plain}
\newtheorem{theorem}[subsection]{Theorem}
\newtheorem{proposition}[subsection]{Proposition}
\newtheorem{lemma}[subsection]{Lemma}

\theoremstyle{definition}
\newtheorem{definition}[subsection]{Definition}
\newtheorem{example}[subsection]{Example}
\newtheorem{exercise}[subsection]{Exercise}
\newtheorem{situation}[subsection]{Situation}

\theoremstyle{remark}
\newtheorem{remark}[subsection]{Remark}
\newtheorem{remarks}[subsection]{Remarks}

\numberwithin{equation}{subsection}

% Macros
%
\def\lim{\mathop{\mathrm{lim}}\nolimits}
\def\colim{\mathop{\mathrm{colim}}\nolimits}
\def\Spec{\mathop{\mathrm{Spec}}}
\def\Hom{\mathop{\mathrm{Hom}}\nolimits}
\def\Ext{\mathop{\mathrm{Ext}}\nolimits}
\def\SheafHom{\mathop{\mathcal{H}\!\mathit{om}}\nolimits}
\def\SheafExt{\mathop{\mathcal{E}\!\mathit{xt}}\nolimits}
\def\Sch{\mathit{Sch}}
\def\Mor{\mathop{\mathrm{Mor}}\nolimits}
\def\Ob{\mathop{\mathrm{Ob}}\nolimits}
\def\Sh{\mathop{\mathit{Sh}}\nolimits}
\def\NL{\mathop{N\!L}\nolimits}
\def\CH{\mathop{\mathrm{CH}}\nolimits}
\def\proetale{{pro\text{-}\acute{e}tale}}
\def\etale{{\acute{e}tale}}
\def\QCoh{\mathit{QCoh}}
\def\Ker{\mathop{\mathrm{Ker}}}
\def\Im{\mathop{\mathrm{Im}}}
\def\Coker{\mathop{\mathrm{Coker}}}
\def\Coim{\mathop{\mathrm{Coim}}}

% Boxtimes
%
\DeclareMathSymbol{\boxtimes}{\mathbin}{AMSa}{"02}

%
% Macros for moduli stacks/spaces
%
\def\QCohstack{\mathcal{QC}\!\mathit{oh}}
\def\Cohstack{\mathcal{C}\!\mathit{oh}}
\def\Spacesstack{\mathcal{S}\!\mathit{paces}}
\def\Quotfunctor{\mathrm{Quot}}
\def\Hilbfunctor{\mathrm{Hilb}}
\def\Curvesstack{\mathcal{C}\!\mathit{urves}}
\def\Polarizedstack{\mathcal{P}\!\mathit{olarized}}
\def\Complexesstack{\mathcal{C}\!\mathit{omplexes}}
% \Pic is the operator that assigns to X its picard group, usage \Pic(X)
% \Picardstack_{X/B} denotes the Picard stack of X over B
% \Picardfunctor_{X/B} denotes the Picard functor of X over B
\def\Pic{\mathop{\mathrm{Pic}}\nolimits}
\def\Picardstack{\mathcal{P}\!\mathit{ic}}
\def\Picardfunctor{\mathrm{Pic}}
\def\Deformationcategory{\mathcal{D}\!\mathit{ef}}
\usepackage{cleveref}
\crefname{originallemma}{original lemma}{original lemmas}
\Crefname{originallemma}{Original lemma}{Original lemmas}
\makeatletter
\newlabel{algebra-lemma-fun}{{05G8}{source}{}{sourcepointer-1}{}}
\newlabel{algebra-lemma-fun@cref}{{[originallemma][0][]05G8}{[1][0][]source}{}{}{}}
\makeatother
\crefname{originallemma}{original lemma}{original lemmas}
\Crefname{originallemma}{Original lemma}{Original lemmas}
\makeatletter
\newlabel{lemma-canonical-element-well-defined}{{36985}{source}{}{sourcepointer-2}{}}
\newlabel{lemma-canonical-element-well-defined@cref}{{[originallemma][0][]36985}{[1][0][]source}{}{}{}}
\makeatother
\crefname{originallemma}{original lemma}{original lemmas}
\Crefname{originallemma}{Original lemma}{Original lemmas}
\makeatletter
\newlabel{lemma-det-ses-functorial}{{36337}{source}{}{sourcepointer-3}{}}
\newlabel{lemma-det-ses-functorial@cref}{{[originallemma][0][]36337}{[1][0][]source}{}{}{}}
\makeatother
\crefname{originallemma}{original lemma}{original lemmas}
\Crefname{originallemma}{Original lemma}{Original lemmas}
\makeatletter
\newlabel{lemma-det-switch}{{36420}{source}{}{sourcepointer-4}{}}
\newlabel{lemma-det-switch@cref}{{[originallemma][0][]36420}{[1][0][]source}{}{}{}}
\makeatother
\title{Surjectivity up to homotopy of two maps between bounded complexes of finite free modules detects the vanishing of the determinant}
\author{The Stacks Project Authors}
\date{stacks-project:more-algebra.tex:local-snapshot-20261012; GNU FDL 1.2 or later}
\begin{document}
\maketitle
\setcounter{subsection}{1}
\noindent\textit{Editorial note.} All delivered mathematical text is the author's own, copied from the pinned Stacks Project chapter file more-algebra.tex; no statement or proof step was rewritten, reordered or removed.\par\medskip
\noindent\textit{Editorial note.} The chapter file pinned-original-source.tex carries only \textbackslash{}input\{preamble\} on its first line, so no preamble text lives inside the chapter. The retained file sources/stacks-more-algebra/retained-preamble.tex is a mechanical, line-for-line copy of the project's single shared preamble as staged in the snapshot; nothing was added, removed or redefined and no package was added.\par\medskip
\noindent\textit{Editorial note.} The source URL is the source-level Stacks Project URL, as required when more than one card is produced from one source; the per-claim tag is carried by claim\_key.\par\medskip
\noindent\textit{Editorial note.} The four references inside the counted proof span are recorded as non-load-bearing navigation pointers: algebra-lemma-fun resolves to a cross-chapter tag, while lemma-canonical-element-well-defined, lemma-det-ses-functorial and lemma-det-switch point to same-chapter lines. The card S7-alg-det-switch counts lemma-det-switch separately at its own statement and proof lines, and that separate counting is not disturbed here.\par\medskip
\noindent\textit{Allowed standard prerequisites (editorial).}\begin{itemize}
\item The determinant functor on bounded complexes of finite free modules and its canonical element. fixed earlier in the same section and invoked through the cited same-chapter lemmas in the counted proof span. Supplies the determinant whose vanishing controls the surjectivity statement.
\item Homotopy of maps between complexes. the homotopy hypothesis in the counted statement span. Is the relation that lets two maps be compared after passing to the determinant.
\end{itemize}\medskip
\ifdefined\hypertarget\hypertarget{sourcepointer-1}{}\fi
\noindent\textit{Source reference.} Original lemma 05G8: focused-seeds/AGP/sources/stacks-snapshot-20261011/tags.txt, source lines 7052--7052. Cross-chapter navigation pointer into the Stacks Project; the standard fact it supplies is applied inside the counted proof span and the argument resting on it stays there.\par\medskip
\ifdefined\hypertarget\hypertarget{sourcepointer-2}{}\fi
\noindent\textit{Source reference.} Original lemma 36985: sources/stacks-more-algebra/pinned-original-source.tex, source lines 36985--36985. Same-chapter navigation pointer; the statement it names is standard for the section and the argument that consumes it stays inside the counted proof span.\par\medskip
\ifdefined\hypertarget\hypertarget{sourcepointer-3}{}\fi
\noindent\textit{Source reference.} Original lemma 36337: sources/stacks-more-algebra/pinned-original-source.tex, source lines 36337--36337. Same-chapter navigation pointer; the statement it names is standard for the section and the argument that consumes it stays inside the counted proof span.\par\medskip
\ifdefined\hypertarget\hypertarget{sourcepointer-4}{}\fi
\noindent\textit{Source reference.} Original lemma 36420: sources/stacks-more-algebra/pinned-original-source.tex, source lines 36420--36420. Same-chapter navigation pointer; the statement it names is standard for the section and the argument that consumes it stays inside the counted proof span.\par\medskip
\begin{lemma}
\label{lemma-homotopic-surjections}
Let $R$ be a ring. Let $a^\bullet : K^\bullet \to L^\bullet$ be a map of
complexes of $R$-modules satisfying (1), (2), (3) above.
Let $h : K^0 \to L^{-1}$ be a map such that
$b^0 = a^0 + d \circ h$ and $b^{-1} = a^{-1} + h \circ d$ are surjective.
Then $\det(a^\bullet) = \det(b^\bullet)$ as maps
$\det(K^\bullet) \to \det(L^\bullet)$.
\end{lemma}

\begin{proof}
Suppose there exists a map $\tilde h : K^0 \to K^{-1}$ such that
$h = a^{-1} \circ \tilde h$ and such that
$k^0 = \text{id} + d \circ \tilde h : K^0 \to K^0$ and
$k^1 = \text{id} + \tilde h \circ d : K^{-1} \to K^{-1}$ are isomorphisms.
Then we obtain a commutative diagram
$$
\xymatrix{
0 \ar[r] &
\Ker(b^\bullet) \ar[r] \ar[d]_{c^\bullet} &
K^\bullet \ar[r]_{b^\bullet}
\ar[d]_{k^\bullet} &
L^\bullet \ar[r] \ar[d]^{\text{id}} &
0 \\
0 \ar[r] &
\Ker(a^\bullet) \ar[r] &
K^\bullet \ar[r]^{a^\bullet} &
L^\bullet \ar[r] &
0
}
$$
of complexes, where $c^\bullet$ is the induced isomorphism of kernels.
Using Lemma \ref{lemma-det-ses-functorial} we see that
$$
\xymatrix{
\det(\Ker(b^i)) \otimes \det(L^i) \ar[r] \ar[d]_{\det(c^i) \otimes 1} &
\det(K^i) \ar[d]^{\det(k^i)} \\
\det(\Ker(a^i)) \otimes \det(L^i) \ar[r] &
\det(K^i)
}
$$
commutes. Since $\det(c^\bullet)$ maps the canonical trivialization
of $\det(\Ker(a^\bullet))$ to the canonical trivializatio of $\Ker(b^\bullet)$
(Lemma \ref{lemma-canonical-element-well-defined})
we see that we conclude if (and only if)
$$
\det(k^0) = \det(k^{-1})
$$
as elements of $R$ which follows from Lemma \ref{lemma-det-switch}.

\medskip\noindent
Suppose there exists a direct summand $U \subset K^{-1}$ such that both
$a^{-1}|_U : U \to L^{-1}$ and $b^{-1}|_U : U \to L^{-1}$
are isomorphisms. Define $\tilde h$ as the composition of $h$
with the inverse of $a^{-1}|_U$. We claim that $\tilde h$ is a
map as in the first paragraph of the proof. Namely, we have
$h = a^{-1} \circ \tilde h$ by construction. 
To show that $k^{-1} : K^{-1} \to K^{-1}$
is an isomorphism it suffices to show that it is surjective
(Algebra, Lemma \ref{algebra-lemma-fun}). Let $u \in U$. We may
choose $u' \in U$ such that $b^{-1}(u') = a^{-1}(u)$.
Then $u = k^{-1}(u')$. Namely, both $u$ and $k^{-1}(u')$
are in $U$ and $a^{-1}(u) = a^{-1}(k^{-1}(u'))$ by a
calculation\footnote{$a^{-1}(k^{-1}(u')) =
a^{-1}(u') + a^{-1}(\tilde h(d(u'))) =
a^{-1}(u') + h(d(u')) = b^{-1}(u') = a^{-1}(u)$} Since $a^{-1}|_U$
is an isomorphism we get the equality. Thus $U \subset \Im(k^{-1})$.
On the other hand, if $x \in \Ker(a^{-1})$ then
$x = k^{-1}(x) \bmod U$. Since $K^{-1} = \Ker(a^{-1}) + U$ we conclude
$k^{-1}$ is surjective.
Finally, we show that $k^0 : K^0 \to K^0$ is surjective.
First, since $a^0 \circ k^0 = b^0$ we see that $a^0 \circ k^0$ is
surjective. If $x \in \Ker(a^0)$, then $x = d(y)$ for some
$y \in \Ker(a^{-1})$. We may write $y = k^{-1}(z)$ for some $z \in K^{-1}$
by the above. Then $x = k^0(d(z))$ and we conclude.

\medskip\noindent
Final step of the proof. It suffices to find $U$ as in the 
preceding paragraph, but this may not always be possible. However,
in order to show equality of two maps of $R$-modules, it suffices
to do so after localization at primes of $R$. Hence we may assume
$R$ is local. Then we get the following problem: suppose
$$
\alpha, \beta : R^{\oplus n} \longrightarrow R^{\oplus m}
$$
are two surjective $R$-linear maps. Find a direct summand
$U \subset R^{\oplus n}$ such that both $\alpha|_U$ and $\beta|_U$
are isomorphisms. If $R$ is a field, this is possible by linear algebra.
In general, one takes a solution over the residue field and lifts this
to a solution over the local ring $R$. Some details omitted.
\end{proof}

\end{document}
